\chapter{The complete tinyhttp framework}
\chapterintro{The pieces now form a small framework: parser, public message types, router, middleware application, and bounded server. This chapter follows one request across every interface.}

\begin{objectives}
\begin{itemize}
  \item Use the public framework API without depending on implementation internals.
  \item Trace data from socket bytes to handler and back to response bytes.
  \item Recognize invariants guaranteed at each boundary.
  \item Run the package tests and extend it without widening its protocol claims silently.
\end{itemize}
\end{objectives}

\section{Package boundaries}

The final \code{tinyhttp} package contains five focused modules:

\begin{center}
\begin{tabularx}{\textwidth}{@{}>{\bfseries}p{28mm} X X@{}}
\toprule
Module & Owns & Does not own \\
\midrule
models & Headers, Request, Response, HTTPError & socket I/O or routing \\
parser & HTTP/1.1 request syntax, framing, target parsing & JSON domain rules \\
router & method/path patterns and parameters & authorization or database lookup \\
application & middleware, dispatch, safe error boundary & TCP lifecycle \\
server & listener, workers, timeouts, keep-alive, serialization & task semantics \\
\bottomrule
\end{tabularx}
\end{center}

\code{tinyhttp.\_\_init\_\_} exports only the teaching API. Application code
imports \code{Application}, \code{Request}, \code{Response}, \code{HTTPError},
\code{Server}, and \code{ServerConfig}; it need not import parser details.

\section{A minimal application}

The public shape is intentionally familiar but transparent:

\begin{lstlisting}[style=pythonbook]
from tinyhttp import Application, Request, Response, Server

app = Application()

@app.route("GET", "/hello/{name}")
def hello(request: Request) -> Response:
    return Response.json(
        {"message": f"Hello, {request.path_params['name']}!"}
    )

Server(app).serve_forever()
\end{lstlisting}

The decorator registers a handler. The handler accepts one Request and returns
one Response. The server depends on the application's \code{handle} method.
There is no global hidden request object and no automatic database injection.

\section{Trace one request}

For \method{GET} \code{/hello/Ada?formal=yes}:

\begin{enumerate}
  \item The listener accepts a connected socket and a worker takes it from the queue.
  \item A connection-local parser receives bytes until one complete message exists.
  \item The parser validates HTTP/1.1, Host, fields, and framing; it decodes the
        path and creates query lists.
  \item \code{Application.handle} composes middleware and calls its dispatcher.
  \item The router matches \code{/hello/\{name\}} and sets
        \code{path\_params["name"]} to \code{"Ada"}.
  \item The handler creates a UTF-8 JSON Response.
  \item Middleware unwinds and can inspect the response.
  \item The server adds wire fields, computes length, serializes, and calls
        \code{sendall}.
  \item If policy permits, the parser and socket remain for the next request.
\end{enumerate}

Each transition narrows the type of uncertainty. Socket code sees bytes; parser
code produces an HTTP request; router code produces parameters; the handler
produces a response value; the server produces bytes again.

\begin{figure}[htbp]
\centering
\begin{tikzpicture}[
  node distance=4mm,
  box/.style={draw=Rule,fill=PaperTint,minimum height=8mm,
              text width=26mm,align=center,font=\small\sffamily},
  arrow/.style={-{Latex[length=2mm]},Accent}
]
\node[box] (socket) {socket bytes};
\node[box,right=of socket] (parser) {RequestParser};
\node[box,right=of parser] (app) {middleware};
\node[box,right=of app] (router) {router};
\node[box,below=8mm of router] (handler) {handler};
\node[box,left=of handler] (response) {Response};
\node[box,left=of response] (serializer) {serializer};
\node[box,left=of serializer] (wire) {socket bytes};
\draw[arrow] (socket) -- (parser);
\draw[arrow] (parser) -- (app);
\draw[arrow] (app) -- (router);
\draw[arrow] (router) -- (handler);
\draw[arrow] (handler) -- (response);
\draw[arrow] (response) -- (serializer);
\draw[arrow] (serializer) -- (wire);
\end{tikzpicture}
\caption{The framework data flow. The return path is represented on the lower row.}
\label{fig:framework-flow}
\end{figure}

\section{Configuration is policy}

\code{ServerConfig} centralizes address, port, workers, queue size, listen backlog,
read timeout, keep-alive, request count, and parser limits. Tests replace the
port with zero and use fewer workers. An application can disable keep-alive
without rewriting server logic.

Configuration values should not be changed independently without thought. A
larger body limit increases memory pressure. More workers increase simultaneous
database contention. A longer timeout multiplies the duration for which slow
clients hold capacity.

\section{Run the verification suite}

From \filepath{code/final}:

\begin{wiremessage}
python -m unittest discover -s tests -v
\end{wiremessage}

The suite tests models and parser, application/router/middleware, real loopback
socket exchange, and the Task API. Integration tests bind an ephemeral port and
stop the server in teardown. A sandbox that forbids local sockets must grant that
specific capability; skipping every network test would leave the central
artifact unverified.

\section{How to extend honestly}

Before adding a feature, update the contract:

\begin{enumerate}
  \item name the semantics and relevant standard rules;
  \item define accepted and rejected inputs;
  \item identify which module owns the feature;
  \item add adversarial tests before widening acceptance;
  \item update documentation and limits;
  \item preserve behavior for existing clients or version the interface.
\end{enumerate}

Adding chunked bodies, for example, is not just decoding hexadecimal sizes. It
affects framing precedence, decoded limits, trailers, connection buffering,
tests, and intermediary compatibility. Honest software names the whole change.

\begin{checkpoint}
You now have a reusable educational framework. Before beginning the backend,
you should be able to open any module and state its inputs, outputs, invariants,
resource ownership, and failure modes. The next part uses the framework rather
than changing HTTP rules casually.
\end{checkpoint}

\summarytitle

\code{tinyhttp} is small because its boundaries are narrow, not because it hides
work. Five modules transform connection bytes into requests, route through
middleware and handlers, and serialize responses under explicit limits. The
public API is enough to build the Task backend while implementation invariants
remain inspectable.

\reviewtitle
\begin{enumerate}
  \item Which module owns each request-to-response transition?
  \item What invariants may a handler assume about its Request?
  \item Why is ServerConfig more than cosmetic settings?
  \item What must change before claiming support for chunked requests?
\end{enumerate}

\begin{exercise}[Add a route safely]
Add \code{GET /echo/\{value\}} to a small application, returning JSON with the
path value and repeated \code{tag} query values. Test Unicode percent encoding,
blank query values, HEAD fallback, OPTIONS, and a wrong method.
\end{exercise}

\begin{exercise}[Boundary review]
Choose one class from each module and write its contract as preconditions,
postconditions, invariants, and possible errors. Compare those contracts with
the actual tests and add one missing case.
\end{exercise}

